\chapter{Social Background of Indian Nationalism}

\section{The Book \& the Mode-of-Production Method}

\emph{Social Background of Indian Nationalism} (A. R. Desai, 1948) analyses the Indian national movement from 1830 to 1947 through the Marxist lens --- i.e.\ through the \textbf{mode of production}, not culture or rituals. Its sequel is \emph{Recent Trends in Indian Nationalism} (1960), and later \emph{State and Society in India} (1975).

\begin{QuickRevision}
\begin{itemize}
\item Book (1948) analyses nationalism 1830--1947 through mode of production.
\item Sequel: Recent Trends in Indian Nationalism (1960); State and Society in India (1975).
\end{itemize}
\end{QuickRevision}

\section{Pre-British Feudal Mode of Production}

Desai's five features of the pre-British feudal mode of production:
\begin{itemize}
\item \textbf{Village as community owner of land} --- no private property.
\item \textbf{Self-sufficiency} --- subsistence economy, barter, no classes.
\item \textbf{Monarch receives produce through revenue collectors} (representative of the monarch; monarch = symbolic owner).
\item \textbf{Static economy}.
\item \textbf{Caste as the steel frame of Hinduism} --- occupation decided by caste (Baniyas were traders, Brahmins were educators).
\end{itemize}

\begin{QuickRevision}
\begin{itemize}
\item Feudal features: communal land, self-sufficiency, revenue collector, static economy, caste as steel frame.
\end{itemize}
\end{QuickRevision}

\section{The Colonial Mode of Production}

With the British, \textbf{means of production changed} (in Europe the industrial revolution came from below; in India it was super-imposed from above --- East India Company, ports, railways, mining). The five features changed:
\begin{itemize}
\item Village commune \textbf{eroded} by peasant proprietors (zamindari, Permanent Settlement, private ownership, sunset laws).
\item Self-sufficiency ended --- villages interconnected (postal, telegraph, railways, factories in Chhota Nagpur, mining, deforestation leading to tribal uprisings for material needs).
\item Monarchy became a \textbf{centralized state} (British India, Supreme Court at Kolkata = all-India, Governor-General).
\item Static became a \textbf{mobile/evolving economy}.
\item Caste as steel frame \textbf{rusting} --- the class of artisans disappeared (modern industry); new \textbf{rural classes} (zamindar, absentee landlords/gentlemen farmers, moneylenders, agricultural labourers, tenants, petty-bourgeois merchants) and \textbf{urban classes} (capitalist industrialists, working class, petty traders, professional/white-collar classes) emerged.
\end{itemize}

\begin{KeyConcept}
This is a \textbf{qualitative structural transformation} --- a change from one mode of production to another, not a quantitative change.
\end{KeyConcept}

\begin{QuickRevision}
\begin{itemize}
\item Means of production changed (imposed from above).
\item Village commune gave way to peasant proprietors; self-sufficiency to interconnection; monarchy to centralized state; static to mobile economy; caste rusting with new rural/urban classes.
\item Qualitative structural transformation.
\end{itemize}
\end{QuickRevision}

\section{Factors Leading to Nationalist Consciousness}

Nationalism is a \textbf{form of consciousness} (ideology in the superstructure, shaped by material conditions --- unlike Hegel, who placed ideology at the base). The trigger factors:
\begin{itemize}
\item \textbf{Introduction of the press} --- Benedict Anderson's \emph{Imagined Communities}; gazettes, pamphlets, journals, exchange of letters.
\item \textbf{Railways} --- enabled Gandhi to travel and lead movements.
\item \textbf{Postal services} --- (they later acquired new functions, e.g.\ postal banking).
\item \textbf{Modern industry}.
\item \textbf{Centralized uniform laws} --- replacing caste panchayats and local laws with all-India laws.
\item \textbf{English education} --- one universal language for communication (about 5\% English literacy).
\end{itemize}

\begin{CriticalPoint}
Nationalism was the \textbf{unanticipated consequence} of modernisation --- never planned by the British.
\end{CriticalPoint}

\begin{QuickRevision}
\begin{itemize}
\item Nationalism = form of consciousness (superstructure, shaped by material conditions).
\item Factors: press (Anderson), railways, postal, industry, uniform laws, English education.
\item Nationalism = unanticipated consequence of modernisation.
\end{itemize}
\end{QuickRevision}

\section{The Five Phases of Indian Nationalism}

Each phase is characterised by a specific class and its interest:
\begin{itemize}
\item \textbf{Phase 1 (1820s--1885)} --- the \textbf{intelligentsia} class; social-religious reforms (Raja Ram Mohan Roy, I. C. Vidyasagar); narrow social base; the \textbf{Indian renaissance} (reawakening); primary westernized people. Reforms: abolition of sati (William Bentinck), child marriage, widow remarriage. Feminism in India was introduced by men (British men + westernized Hindu men), unlike the West (Mary Wollstonecraft).
\item \textbf{Phase 2 (1885--1905)} --- INC established (all-India), Swadeshi movement (1905); social base widened; three classes --- \textbf{educated middle class} (modern schools), \textbf{merchant class} (Indian international trade), \textbf{industrialist class} (modern industry), each a specific caste group (Brahmins, banias, Parsis). Congress under moderates served these classes. Fissure leading to the \textbf{Surat Split (1907)} = polarization (moderates = upper classes, extremists = lower classes).
\item \textbf{Phase 3 (1905--1918)} --- home rule; Gandhi entered politics with mass-line politics; \textbf{militantization/Hinduization} of the movement (Shivaji, Ganga, Ganpati --- Hindu symbols alienated Muslims; nationalism equated with Hinduism).
\item \textbf{Phase 4 (till 1934)} --- civil disobedience (1930--34); social base enormously enlarged (Ahmedabad mill strike 1918, Champaran, Kheda); class to mass movement. WWI, the Russian revolution, the Khilafat movement (unified Hindus-Muslims) stirred the masses. Congress remained under the influence of the capitalist class (swadeshi/boycott served industrialists who financed Congress); Gandhi = class collaborationist.
\item \textbf{Phase 5 (1934--1939)} --- disenchantment with Gandhian ideology; \textbf{Congress Socialist Party} (workers and peasants on class lines); \textbf{All India Kisan Sabha} (pressure group of the CPI; objectives: abolition of landlordism, repudiation of debts, socialist state); growth of the Communist Party; Hindustan Socialist Republican Association (Bhagat Singh).
\end{itemize}

\begin{QuickRevision}
\begin{itemize}
\item Five phases, each = a class and its interest: intelligentsia (reforms), three classes (INC/swadeshi), Hinduization, mass movement, socialist turn.
\item Surat Split 1907 = polarization; CSP and Kisan Sabha in phase 5.
\end{itemize}
\end{QuickRevision}

\section{Conclusion: External vs Internal Colonialism}

By 1946 (after Quit India 1942) the social base had enormously enlarged, but Congress remained dominated by the dominant class. Desai concludes: \textbf{external colonialism was replaced by internal colonialism} --- the British left, and the Congress (representing the Indian capitalist class) took over. The nationalist movement was an \textbf{all-class movement}, yet it remained dominated by the Gandhian Congress, which then formed the Indian state and chose \textbf{bourgeois (not socialist) industrialization}.

\begin{CriticalPoint}
Desai's contention with Marx: Marx thought capitalism would revolutionise the mode of production; Desai held that capitalism in India \textbf{destroyed the indigenous industry} that could have developed capitalism on its own. Desai was a nationalist as well as a Marxist --- more \textbf{Trotskyite} than Marxist (against the two-stage revolution idea).
\end{CriticalPoint}

\begin{QuickRevision}
\begin{itemize}
\item External colonialism replaced by internal colonialism; Indian state chose bourgeois industrialization.
\item Desai: capitalism destroyed indigenous industry; more Trotskyite than Marxist.
\end{itemize}
\end{QuickRevision}
